% Compile with: xelatex seosong_result_report.tex
\documentclass[11pt]{article}

\usepackage[a4paper,margin=1in]{geometry}
\usepackage{fontspec}
\usepackage{xeCJK}
\usepackage{booktabs}
\usepackage{graphicx}
\usepackage{hyperref}
\usepackage{siunitx}
\usepackage{array}
\usepackage{caption}
\usepackage{float}
\usepackage{xcolor}

\setmainfont{Times New Roman}
\setsansfont{Arial}
\setCJKmainfont{Apple SD Gothic Neo}
\setmonofont{Menlo}
\setCJKmonofont{Apple SD Gothic Neo}
\hypersetup{
  colorlinks=true,
  linkcolor=blue!50!black,
  urlcolor=blue!50!black
}
\sisetup{
  round-mode=places,
  round-precision=4,
  table-number-alignment=center
}

\title{\textbf{서송병원 Forecasting Result Report}\\
\large LAVAR Seed Stability, Blend Robustness, and Ensemble Interval Results}
\author{LAVAR Health-Economic Forecasting Study}
\date{June 2026}

\begin{document}
\maketitle

\begin{abstract}
This report summarizes the 서송병원 forecasting results produced by the
Seosong replication workflow. The evaluation uses daily patient features
and prescription targets from 2020-12-14 through 2025-12-31, with 1,844
daily observations, 38 patient-side covariates, and 3,126 prescription
targets. Across five random seeds and 79 rolling folds per seed, the
S11 latent-16 configuration improves average MSE over the S04 GRU
configuration, but with materially higher seed-to-seed variance and tail
risk. A fixed ensemble, $0.55\times$S11 $+0.45\times$S04, further reduces
mean MSE relative to pure S11 by 0.0596, or 3.17\%, but does not satisfy
the promotion rule because p95 fold MSE worsens relative to pure S11.
\end{abstract}

\section{Data and Evaluation Setup}

The source artifacts are the Seosong report outputs from the seed
stability, blend robustness, and ensemble workflows. The base data files are:
\begin{itemize}
  \item \texttt{data\_patient.parquet}: 1,844 rows by 38 columns.
  \item \texttt{data\_prescription.parquet}: 1,844 rows by 3,126 columns.
  \item Date range: 2020-12-14 to 2025-12-31.
\end{itemize}

Each scenario was evaluated over five seeds:
$\{42, 123, 456, 789, 1024\}$. The rolling evaluation used a 14-day
forecast horizon, 14-day fold step, 90-day retraining cadence, quality
triggers enabled, and 79 folds per run.

\section{Model Scenarios}

The Seosong workflow compares two base scenarios and one fixed blend:
\begin{itemize}
  \item \textbf{S04\_GRU\_H32\_L1}: GRU temporal decoder baseline.
  \item \textbf{S11\_LATENT16}: latent dimension 16 configuration.
  \item \textbf{Blend w=0.55}: $0.55\times$S11 $+0.45\times$S04, clipped
  to non-negative predictions.
\end{itemize}

The ensemble interval artifact computes empirical 95\% intervals using
the five seeded blend predictions.

\section{Seed Stability Results}

Table~\ref{tab:seed-stability} reports seed-level stability summaries from
\texttt{seed\_stability\_summary.csv}. S11 has the better mean MSE and MAE,
but its coefficient of variation for mean MSE is roughly six times larger
than S04's, indicating a less stable seed response.

\begin{table}[H]
\centering
\caption{Seed stability by base scenario.}
\label{tab:seed-stability}
\begin{tabular}{lrrrrrr}
\toprule
Scenario & Seeds & Mean MSE & SD MSE & CV MSE & Mean MAE & p95 Fold MSE \\
\midrule
S04 GRU H32 L1 & 5 & 2.0234 & 0.0540 & 0.0267 & 0.2712 & 4.6293 \\
S11 LATENT16 & 5 & 1.8810 & 0.2986 & 0.1587 & 0.2616 & 4.4179 \\
\bottomrule
\end{tabular}
\end{table}

The S11 mean MSE improvement over S04 is 0.1424, equal to a 7.04\%
reduction relative to S04. However, S11's maximum fold MSE reaches
20.8928 in its worst seed, compared with 13.7482 for S04.

\section{Blend and Ensemble Results}

Table~\ref{tab:ensemble-overall} summarizes the fixed blend against pure
S04 and pure S11 from the ensemble metrics table. The blend is the best
option by average MSE and average maximum-fold MSE, but pure S11 remains
best by p95 fold MSE.

\begin{table}[H]
\centering
\caption{Overall ensemble comparison.}
\label{tab:ensemble-overall}
\begin{tabular}{lrrrrr}
\toprule
Label & Seeds & Mean MSE & SD MSE & Mean MAE & p95 Fold MSE \\
\midrule
Blend w=0.55 & 5 & 1.8214 & 0.1641 & 0.2615 & 4.4676 \\
Pure S04 & 5 & 2.0210 & 0.0520 & 0.2708 & 4.6053 \\
Pure S11 & 5 & 1.8810 & 0.2986 & 0.2616 & 4.4179 \\
\bottomrule
\end{tabular}
\end{table}

Relative to pure S11, the fixed blend reduces mean MSE by 0.0596
($-3.17\%$) and reduces the seed-level standard deviation of mean MSE
from 0.2986 to 0.1641. Relative to pure S04, the blend reduces mean MSE
by 0.1995 ($-9.87\%$). The tradeoff is tail performance: blend p95 fold
MSE is 4.4676, which is 0.0497 higher than pure S11.

\section{Robustness Decision}

The fixed-weight robustness decision table reports that both $w=0.50$ and
$w=0.55$ win on mean MSE for three of five seeds, win on p95 fold MSE for
two of five seeds, and win both criteria for only two of five seeds.
Accordingly, \texttt{promote\_blend\_rule=False} for both weights.

\begin{table}[H]
\centering
\caption{Blend promotion decision.}
\label{tab:blend-decision}
\begin{tabular}{rrrrrr}
\toprule
$w$ & Seeds & Mean-MSE Wins & p95 Wins & Both Wins & Mean $\Delta$MSE \\
\midrule
0.50 & 5 & 3 & 2 & 2 & -0.0538 \\
0.55 & 5 & 3 & 2 & 2 & -0.0596 \\
\bottomrule
\end{tabular}
\end{table}

This means the blend is useful as an average-error reducer and variance
stabilizer, but it should not replace pure S11 as the default operating
rule if tail-fold robustness is the governing criterion.

\section{Interval Output}

The final ensemble workflow generates per-drug empirical 95\% intervals
from seeded blend predictions:
\texttt{ensemble\_ci\_per\_drug.csv}. The visualization in
Figure~\ref{fig:true-pred-ci} plots true values, ensemble predictions,
and the empirical interval for selected drug targets.

\begin{figure}[H]
\centering
\includegraphics[width=\textwidth]{ensemble_true_vs_pred_ci_per_drug.png}
\caption{Seosong ensemble true vs. predicted values with empirical 95\% intervals.}
\label{fig:true-pred-ci}
\end{figure}

\section{Conclusion}

For 서송병원, S11 LATENT16 is the stronger single model by average error,
while S04 is more seed-stable. The $w=0.55$ blend gives the best average
MSE and substantially reduces S11 seed variance, but it worsens p95 fold
MSE against pure S11 and therefore fails the current promotion rule. The
recommended reporting stance is:
\begin{itemize}
  \item use S11 as the primary single-model benchmark;
  \item retain the $w=0.55$ blend as a variance-stabilized sensitivity
  analysis and interval-production model;
  \item avoid promoting the blend as the default model until tail-fold
  robustness improves.
\end{itemize}

\end{document}
